\documentclass[11pt,letterpaper]{article}

\usepackage[margin=1in]{geometry}
\usepackage{amsmath,amssymb,amsthm}
\usepackage{mathtools}
\usepackage{mathrsfs}
\usepackage{natbib}
\usepackage{url}
\usepackage{cancel}
\newcommand{\doi}[1]{doi: #1}

% ---- macros (match proof_exposition_v22) ----
\DeclareMathOperator{\Var}{Var}
\DeclareMathOperator{\Cov}{Cov}
\newcommand{\E}{\mathbb{E}}
\newcommand{\R}{\mathbb{R}}
\newcommand{\bX}{\boldsymbol{X}}

% ---- theorem environments ----
\newtheorem{proposition}{Proposition}
\newtheorem{corollary}{Corollary}
\newtheorem{lemma}{Lemma}
\theoremstyle{remark}
\newtheorem*{remark}{Remark}

\title{Quantitative Bounds on Sobol Index Shifts\\
Under Smooth Output Transformations}
\author{}
\date{}

\begin{document}

\maketitle

% ======================================================================
\section{Background and purpose}
\label{sec:background}
% ======================================================================

The companion commutativity memo proves a qualitative statement: a continuous
nonconstant pointwise transform $\phi\colon\R\to\R$ commutes with Sobol index
computation for every admissible scalar model if and only if $\phi$ is affine.
The necessity direction uses a binary midpoint witness model: for each endpoint
pair, the witness forces midpoint preservation of $\phi$, and the
Roberts--Varberg midpoint theorem then converts midpoint preservation plus
continuity into global affinity \citep[Theorem~4.1]{RobertsVarberg1973}. The
witness model is not intended to represent all admissible models; it is a
diagnostic submodel that establishes the midpoint property for an arbitrary
endpoint pair.

This memo asks the next quantitative question. If $Z=\phi(Y)$ is a smooth
output transformation, how large can the shift in Sobol ANOVA component,
closed, or total-effect indices be? A theorem that compares $Z$ only to $Y$
would miss a central mechanism: when the first derivative of $\phi$ vanishes at
the mean of $Y$, the transformed output may be governed by a quadratic or
higher-order Taylor term rather than by an affine rescaling of $Y$. For example,
if $Y=X_1+X_2$ and $\phi(y)=y^2$ at a centered mean, then
$\phi(Y)=X_1^2+2X_1X_2+X_2^2$ introduces an interaction term into an otherwise
additive model.

The main result therefore uses a unified perturbation formulation. Instead of
forcing the reference output to be $Y$, Proposition~\ref{prop:abstract_bound}
compares Sobol indices of two arbitrary square-integrable outputs $Z=V+E$ when
the centered perturbation $E-\E E$ is small in $L^2$ relative to the standard
deviation of the reference $V$. Corollary~\ref{cor:taylor_bound} then applies
this abstract result with the Taylor-polynomial reference
\begin{equation}
  V_m=\sum_{k=1}^{m}\frac{\phi^{(k)}(\mu_Y)}{k!}(Y-\mu_Y)^k.
\end{equation}
The nonzero-slope local-affine bound is the special case $m=1$ with
$\phi'(\mu_Y)\neq0$. The critical-point case is not excluded: if the first
nonzero Taylor term is quadratic and the resulting quadratic reference is
nonconstant, choose $m=2$ and compare $\phi(Y)$ to that reference. For an exactly
polynomial transform of degree at most $m$, the Taylor residual is zero and the
comparison is exact whenever the polynomial reference has positive variance.

No new sensitivity measure is introduced. The only additional notation below is
a normalized-projection shorthand used to prove bounds simultaneously for Sobol
ANOVA component indices, closed indices, and total-effect indices.

% ======================================================================
\section{Literature review}
\label{sec:lit}
% ======================================================================

\subsection{Summary of the companion memo}

The companion memo establishes the following context, summarized here for
completeness. The product-measure structure guarantees that a square-integrable
scalar output $Y=f(\bX)$ with independent inputs admits a unique
Hoeffding--Sobol functional ANOVA decomposition into mutually orthogonal
$L^2$-components \citep{Hoeffding1948,Sobol1993}.\footnote{The 1948 Hoeffding
paper introduced U-statistics and a related orthogonality property for symmetric
kernels of two variables; it does not state the variance-based ANOVA decomposition
of a function of $d$ independent inputs. The latter is due to \citet{Sobol1993}.
The joint attribution follows standard practice in the sensitivity analysis
literature \citep[e.g.,][]{Saltelli2008}.}
The variance decomposes as
$\Var[Y]=\sum_{\emptyset\neq u}D_u^Y$, where $D_u^Y=\Var[f_u(\bX_u)]$.
The corresponding Sobol ANOVA component index is
$S_u^Y=D_u^Y/\Var[Y]$ for each nonempty ANOVA subset $u$.
The first-order Sobol index for input $X_i$ is \citep{Sobol1993}
\begin{equation}
  S_i^Y=\frac{D_i^Y}{\Var[Y]},\qquad D_i^Y=\Var_{X_i}\{\E[Y\mid X_i]\},
  \label{eq:sobol_def}
\end{equation}
where the conditional expectation integrates over all randomness except $X_i$
under the product input law. The total-effect index is
$(S_i^T)^Y=(\sum_{u\ni i}D_u^Y)/\Var[Y]$ \citep{Homma1996}. Widely used
Monte Carlo designs and estimators for these variance components are discussed
by \citet{Saltelli2002}, \citet{Jansen1999}, \citet{Saltelli2008}, and
\citet{Saltelli2010}. For the independent-input variance game,
\citet{Owen2014} shows that first-order and total-effect Sobol indices bracket
the corresponding Shapley effect.

Affine output transformations $Z=aY+b$ ($a\neq 0$) preserve all Sobol indices
because every ANOVA component and the total variance are multiplied by the same
factor $a^2$. For Hilbert-space-valued outputs, \citet{Gamboa2014} prove that
generalized Sobol indices are invariant under isometries, scalings, and
translations; in the scalar case this reduces to the same affine cancellation.
The companion memo proves that affine maps are the \emph{only} continuous
nonconstant pointwise transforms that commute with Sobol index computation for
every admissible scalar model. That result is qualitative: it guarantees a
shift exists for any non-affine $\phi$, but gives no bound on its magnitude.
The present memo supplies a complementary quantitative bound. It should not be
read as a second commutativity theorem: affine maps remain the only universal
commuting transformations. The bound here explains how far the Sobol structure of
$\phi(Y)$ can lie from the Sobol structure of an appropriate reference output,
usually a Taylor polynomial in $Y-\E Y$.

\subsection{Transformation invariance and its limits}

\citet{Borgonovo2014} study which global sensitivity statistics retain their
interpretation under monotone output transformations. They show that classical
variance-based indices, including Sobol indices, are not invariant under
nonlinear monotone transforms, whereas certain density- and CDF-based statistics
are. The review of \citet{Borgonovo2016} further emphasizes that different
sensitivity measures answer different transformation-dependent questions. For
independent inputs, one-to-one marginal reparametrisations of \emph{inputs}
preserve $S_i^Y$ because the conditioning sigma-fields are unchanged; output
transformations are fundamentally different because they replace $\E[Y\mid X_i]$
with $\E[\phi(Y)\mid X_i]$.
Relatedly, \citet{BaucellsBorgonovo2013} develop invariant probabilistic
sensitivity analysis in a decision-analytic setting by measuring distances
between risk profiles in a way that is invariant to monotone transformations.
This reinforces the main distinction: transformation-invariant measures answer a
distributional or decision-analytic question, whereas Sobol indices allocate
variance on the output scale being analyzed.

Rank transformation of the output is a widely used diagnostic tool.
\citet{Iman1988} and \citet{Saltelli1995} show that recomputing sensitivity
indices after replacing $Y$ by ranks or normal scores yields sensitivity indices
for the transformed response, not invariant versions of the original Sobol
indices. Such rank-based analyses can change the apparent balance between main
and interaction effects; the direction and size of that change are model
specific. This memo studies the analogous question for a smooth deterministic
transform: the output of the calculation is an index for $\phi(Y)$, and the
question is how far it lies from the corresponding index of a chosen reference
output.

\subsection{Derivative-based bounds and distribution-based alternatives}

The closest structural analogue to the present result is the derivative-based
global sensitivity measure (DGSM) literature. \citet{SobolKucherenko2009}
introduced derivative-based global sensitivity measures and showed that
derivative-based quantities can control total-effect Sobol indices in standard
settings. Subsequent work expresses these links through one-dimensional
Poincar\'{e} inequalities: under such inequalities the bound takes the
dimensionless form
\begin{equation}
  (S_i^T)^Y\leq \frac{C_i\,\nu_i}{\Var[Y]},\qquad
  \nu_i=\E\!\left[\!\left(\tfrac{\partial f}{\partial x_i}\right)^{\!2}\right],
  \label{eq:dgsm_bound}
\end{equation}
where $C_i$ is the Poincar\'{e} constant of the marginal distribution of $X_i$.
\citet{Lamboni2013} give general links between DGSMs and total Sobol indices,
and \citet{Roustant2017} develop interval Poincar\'{e} inequalities for applying
these bounds to broader marginal distributions. \citet{Roustant2020} provide
Parseval-type lower bounds on variance-based indices via spectral expansions.
All of these results bound the index itself in terms of derivatives of the model
function $f$. Proposition~\ref{prop:abstract_bound} is complementary: it bounds
index displacement under an output perturbation, and
Corollary~\ref{cor:taylor_bound} expresses that perturbation in terms of Taylor
remainders of the output transform $\phi$.

Moment-independent and distribution-based measures offer an alternative
perspective. Borgonovo's $\delta$ index \citep{Borgonovo2007} measures the
expected $L^1$ distance between the unconditional and conditional output
densities,
\begin{equation}
  \delta_i=\tfrac12\,\E_{X_i}\!\left[
    \int_{\R}\left|p_Y(y)-p_{Y\mid X_i}(y\mid X_i)\right|dy
  \right],
  \label{eq:delta_index}
\end{equation}
and is invariant under strictly monotone output transformations under the usual
regular density-transformation conditions \citep{Borgonovo2014}, making it
immune to the Sobol-index displacement quantified here. Estimation of
$\delta_i$ and related CDF-based measures from given data is studied by
\citet{Plischke2013}. Cram\'{e}r--von Mises sensitivity indices provide another
distribution-oriented alternative: \citet{GamboaKleinLagnoux2018} define indices
based on distances between unconditional and conditional output distributions,
thereby targeting whole-distribution sensitivity rather than variance allocation
alone. Dependence-measure sensitivity indices, including HSIC-based indices,
measure input--output dependence through kernel embeddings \citep{DaVeiga2015}.
They are valuable alternatives to variance allocation, especially for screening,
but fixed-kernel HSIC values generally depend on the chosen kernels and on the
scale or transformation of the output; they should not be described as
Sobol-like invariant variance shares.

Perturbed-law indices study a different perturbation direction: they perturb the
input distribution and track the change in a reliability probability or quantile.
\citet{Lemaitre2015} develop density-modification-based reliability sensitivity,
and \citet{SueurIoossDelage2017} discuss quantile-oriented perturbed-law indices.
Both frameworks ask how a sensitivity-relevant quantity moves under a controlled
perturbation; the perturbation here is output-side transformation rather than
input-law modification.

\subsection{Gap addressed here}

The existing output-transformation literature primarily identifies whether a
sensitivity measure is invariant under a given class of transforms, or recommends
a measure whose target is invariant. Distributional alternatives such as
$\delta$ indices, invariant probabilistic sensitivity analysis, and
Cram\'{e}r--von Mises indices answer different inferential questions. The DGSM
literature supplies bounds on total-effect indices $(S_i^T)^Y$ through
derivatives of the model function $f$; those results do not bound the
displacement caused by replacing one scalar output by another. We have not found
a directly analogous nonasymptotic population bound for the displacement of
Sobol ANOVA component indices or for the standard closed and total-effect sums
of those components under a smooth output transformation. The result below fills
that gap by treating Sobol indices as squared norms of orthogonal ANOVA
projections. The smooth-transform statement is then a Taylor-reference
corollary, which includes both the local-affine and critical-point regimes.

% ======================================================================
\section{Quantitative bound}
\label{sec:bound}
% ======================================================================

This section uses Hilbert-space notation only to make the Sobol algebra
compact. All random variables below live in the same probability space generated
by the independent inputs $X_1,\ldots,X_d$. The space $L^2$ is the set of
square-integrable scalar random variables. Its inner product and norm are
\begin{equation}
  \langle A,B\rangle=\E[AB],
  \qquad
  \|A\|_2=(\E[A^2])^{1/2}.
\end{equation}
After centering a random variable, this norm is its standard deviation:
$\|W-\E W\|_2=\operatorname{sd}(W)$.

Let
\begin{equation}
  \mathcal U_d={u\subseteq{1,\ldots,d}:u\neq\emptyset}
\end{equation}
denote the set of nonempty ANOVA subsets. For each $u\in\mathcal U_d$, let
$\mathcal H_u$ denote the Hoeffding--Sobol ANOVA subspace containing the part of
an output that depends on $\bX_u=(X_j)_{j\in u}$ and is orthogonal to all lower
order components. Under the product input law, these subspaces are mutually
orthogonal, so a centered square-integrable output decomposes as an orthogonal
sum of its ANOVA components. If $\mathscr C\subseteq\mathcal U_d$ is a nonempty
collection of such subsets, define
\begin{equation}
  \mathcal H_{\mathscr C}=\bigoplus_{u\in\mathscr C}\mathcal H_u,
\end{equation}
and let $\Pi_{\mathscr C}$ be the orthogonal projection onto
$\mathcal H_{\mathscr C}$. Intuitively, $\Pi_{\mathscr C}(W-\E W)$ extracts only
the part of the centered output whose ANOVA components belong to the collection
$\mathscr C$.

For any nonconstant square-integrable scalar output $W$, define
\begin{equation}
  W^\circ=W-\E W,
  \qquad
  \mathsf{PE}_{\mathscr C}(W)
  =\frac{\|\Pi_{\mathscr C} W^\circ\|_2^2}{\Var[W]},
  \label{eq:pe_def}
\end{equation}
where $\mathsf{PE}_{\mathscr C}$ means normalized projection energy. This is
proof bookkeeping, not a new sensitivity measure. For $\mathscr C={u}$ it is
the Sobol ANOVA component index $S_u^W$; for the collections in
Corollary~\ref{cor:special_cases} it is the standard closed or total-effect
Sobol index. The theorem is stated for $\mathsf{PE}_{\mathscr C}$ only so that
one projection calculation covers all three standard cases.

\begin{proposition}[Projection perturbation bound for Sobol indices]
\label{prop:abstract_bound}
Let $V,E\in L^2$ be scalar outputs on the same independent-input product space,
and set
\begin{equation}
  Z=V+E,
  \qquad
  \sigma_V^2=\Var[V]>0,
  \qquad
  \eta=\frac{\|E-\E E\|_2}{\sigma_V}.
  \label{eq:eta_def}
\end{equation}
If $\eta<1$, then $0<\Var[Z]<\infty$ and, for every nonempty
$\mathscr C\subseteq\mathcal U_d$,
\begin{equation}
  \left|\mathsf{PE}_{\mathscr C}(Z)-\mathsf{PE}_{\mathscr C}(V)\right|
  \leq
  \frac{
    2\eta\sqrt{p_{\mathscr C}}(1+\sqrt{p_{\mathscr C}})
    +\eta^2(1+p_{\mathscr C})
  }{(1-\eta)^2},
  \qquad
  p_{\mathscr C}=\mathsf{PE}_{\mathscr C}(V).
  \label{eq:abstract_bound}
\end{equation}
Equivalently, the right side may be replaced by its minimum with one, since
Sobol component, closed, and total-effect indices all lie in $[0,1]$.
\end{proposition}

\begin{corollary}[Taylor-reference bound for smooth output transformations]
\label{cor:taylor_bound}
Let $Y=f(\bX)$, where $\bX=(X_1,\ldots,X_d)$ has independent components. Assume
$Y\in[y_-,y_+]$ almost surely, and let $Z=\phi(Y)$ with
$\phi\in C^{m+1}([y_-,y_+])$ for some integer $m\geq1$. Define
\begin{equation}
  \mu_Y=\E[Y],
  \qquad
  U=Y-\mu_Y,
  \qquad
  V_m=\sum_{k=1}^{m}\frac{\phi^{(k)}(\mu_Y)}{k!}U^k,
  \label{eq:taylor_reference}
\end{equation}
and let
\begin{equation}
  R_m=\phi(Y)-\phi(\mu_Y)-V_m.
  \label{eq:taylor_remainder_def}
\end{equation}
Assume $\sigma_m^2:=\Var[V_m]>0$. This is a substantive reference-variance
condition, not merely a statement that some Taylor coefficient is nonzero. A
nonzero polynomial in $U$ can still be constant on the support of $Y$; for
example, if $U^2$ is almost surely constant, a purely quadratic reference has
zero variance and Sobol indices for that reference are undefined. In practice
one should choose $m$ so that the Taylor reference $V_m$ is nonconstant on the
support of the output distribution. If every useful Taylor reference is constant
on that support, the transformed output may itself be constant or may require a
different expansion center or reference construction.
Define
\begin{equation}
  \eta_m=\frac{\|R_m-\E R_m\|_2}{\sigma_m}.
  \label{eq:eta_m_def}
\end{equation}
If $\eta_m<1$, then for every nonempty $\mathscr C\subseteq\mathcal U_d$,
\begin{equation}
  \left|\mathsf{PE}_{\mathscr C}(Z)-\mathsf{PE}_{\mathscr C}(V_m)\right|
  \leq
  \frac{
    2\eta_m\sqrt{p_{m,\mathscr C}}(1+\sqrt{p_{m,\mathscr C}})
    +\eta_m^2(1+p_{m,\mathscr C})
  }{(1-\eta_m)^2},
  \qquad
  p_{m,\mathscr C}=\mathsf{PE}_{\mathscr C}(V_m).
  \label{eq:taylor_bound}
\end{equation}
Moreover, if
\begin{equation}
  \rho_{m+1}=\|\phi^{(m+1)}\|_{\infty,[y_-,y_+]},
\end{equation}
then Taylor's theorem gives the sufficient computable bound
\begin{equation}
  \eta_m
  \leq
  \bar\eta_m
  :=
  \frac{\rho_{m+1}}{(m+1)!\,\sigma_m}
  \left\{\E|Y-\mu_Y|^{2m+2}\right\}^{1/2}.
  \label{eq:eta_m_sufficient}
\end{equation}
Because the right side of \eqref{eq:taylor_bound} is increasing in $\eta_m$ on
$0\leq\eta_m<1$, the same bound remains valid with $\eta_m$ replaced by
$\bar\eta_m$ whenever $\bar\eta_m<1$.
\end{corollary}

\begin{corollary}[Nonzero-slope local-affine case]
\label{cor:local_affine}
Under the assumptions of Corollary~\ref{cor:taylor_bound} with $m=1$, suppose
$c=\phi'(\mu_Y)\neq0$. Then $V_1=c(Y-\mu_Y)$, so
$\mathsf{PE}_{\mathscr C}(V_1)=\mathsf{PE}_{\mathscr C}(Y)$ for every
$\mathscr C$. Define
\begin{equation}
  \sigma_Y^2=\Var[Y],\qquad
  \mu_4=\E[(Y-\mu_Y)^4],\qquad
  \rho_2=\|\phi''\|_{\infty,[y_-,y_+]},
\end{equation}
and
\begin{equation}
  \kappa=\frac{\rho_2\sigma_Y}{|c|},\qquad
  K=\frac{\mu_4^{1/2}}{\sigma_Y^2},\qquad
  \lambda=K\kappa.
  \label{eq:kappa_lambda}
\end{equation}
Then $\eta_1\leq\lambda/2$. Because the right side of
\eqref{eq:abstract_bound} is increasing in $\eta$ on $0\leq\eta<1$, if
$\lambda<2$, then for every nonempty $\mathscr C\subseteq\mathcal U_d$,
\begin{equation}
  \left|\mathsf{PE}_{\mathscr C}(Z)-\mathsf{PE}_{\mathscr C}(Y)\right|
  \leq
  \frac{
    \lambda\sqrt{p_{\mathscr C}}(1+\sqrt{p_{\mathscr C}})
    +\tfrac14\lambda^2(1+p_{\mathscr C})
  }{(1-\lambda/2)^2},
  \qquad
  p_{\mathscr C}=\mathsf{PE}_{\mathscr C}(Y).
  \label{eq:local_affine_bound}
\end{equation}
\end{corollary}

\begin{corollary}[Standard Sobol-index special cases]
\label{cor:special_cases}
The projection notation in Proposition~\ref{prop:abstract_bound} and
Corollaries~\ref{cor:taylor_bound}--\ref{cor:local_affine} gives the following
standard Sobol indices without introducing any new sensitivity measure:
\begin{enumerate}
\item For a nonempty ANOVA subset $u\in\mathcal U_d$, choose
  $\mathscr C=\{u\}$. Then $\mathsf{PE}_{\{u\}}(W)=S_u^W$.
\item For any nonempty input group $A\subseteq\{1,\ldots,d\}$, the closed index is
  \begin{equation}
    (S_A^{\mathrm{cl}})^W
      =\sum_{\emptyset\neq u\subseteq A}S_u^W
      =\frac{\Var\{\E[W\mid\bX_A]\}}{\Var[W]},
  \end{equation}
  obtained by choosing
  $\mathscr C=\{u:\emptyset\neq u\subseteq A\}$.
\item For any nonempty input group $A\subseteq\{1,\ldots,d\}$, the total-effect
  index is
  \begin{equation}
    (S_A^T)^W
      =\sum_{u:u\cap A\neq\emptyset}S_u^W
      =1-\frac{\Var\{\E[W\mid\bX_{-A}]\}}{\Var[W]},
  \end{equation}
  obtained by choosing $\mathscr C=\{u:u\cap A\neq\emptyset\}$. For a single
  input, $(S_i^T)^W=(S_{\{i\}}^T)^W=\sum_{u\ni i}S_u^W$.
\end{enumerate}
Here $\bX_A=(X_j)_{j\in A}$ and $\bX_{-A}=(X_j)_{j\notin A}$. The closed and
total-effect cases are handled as single orthogonal projections onto direct
sums of ANOVA subspaces; no summation of componentwise error bounds and no
combinatorial union-bound factor is required.
\end{corollary}

\begin{remark}[How to read the theorem]
The theorem controls $\phi(Y)$ relative to the reference output $V_m$, not
necessarily relative to $Y$ itself. Three cases are worth distinguishing.

\emph{Local-affine case} ($m=1$, $c=\phi'(\mu_Y)\neq0$): $V_1=c(Y-\mu_Y)$ is
an affine rescaling of $Y$, so $\mathsf{PE}_{\mathscr C}(V_1)=\mathsf{PE}_{\mathscr C}(Y)$
for every $\mathscr C$. The bound in Corollary~\ref{cor:local_affine} is
therefore a direct bound on $|S_u^Z-S_u^Y|$, $|(S_A^{\mathrm{cl}})^Z-(S_A^{\mathrm{cl}})^Y|$,
and $|(S_A^T)^Z-(S_A^T)^Y|$. The condition $\lambda<2$ is a sufficient
stability condition inherited from the abstract requirement $\eta_1<1$. The
bound is still conservative, and it is practically informative only when the
right side of \eqref{eq:local_affine_bound} is substantially smaller than the
trivial cap of one.

\emph{Critical-point and higher-order cases} ($m\geq2$): the reference $V_m$
may contain new ANOVA interactions created by the nonlinear Taylor terms. The
theorem bounds $|S_u^Z-S_u^{V_m}|$, not $|S_u^Z-S_u^Y|$. Large shifts
relative to $S_u^Y$ can occur---for example, squaring an additive model can
create dominant interaction terms---and no uniformly small bound relative to
$Y$ is possible without additional assumptions on the ANOVA structure of the
polynomial reference. The theorem provides the correct quantitative control:
the Sobol indices of $\phi(Y)$ are close to those of the polynomial reference
$V_m$, and the discrepancy from $Y$'s indices is accounted for by the ANOVA
structure of $V_m$ itself.

\emph{Exact polynomial case}: if $\phi$ is a polynomial of degree at most $m$,
then $R_m=0$, $\eta_m=0$, and the theorem says the Sobol indices of $\phi(Y)$
equal those of $V_m$ exactly.
\end{remark}

\begin{remark}[Trivial cap]
Since any component, closed, or total-effect Sobol index lies in $[0,1]$, the
absolute displacement is always at most one. Thus the right side of
\eqref{eq:abstract_bound}, \eqref{eq:taylor_bound}, or
\eqref{eq:local_affine_bound} may always be replaced by its minimum with one.
This cap is mathematically trivial, but it is useful in applications when the
sufficient denominator-stability condition is barely satisfied and the displayed
perturbation bound is consequently conservative.
\end{remark}

\begin{remark}[Population quantities and plug-in estimation]
The inequalities are population statements. In applications, the quantities in
$V_m$, $\sigma_m^2$, $p_{m,\mathscr C}$, and $\eta_m$ may be estimated from a
sensitivity-analysis design, but the displayed inequalities do not include
finite-sample uncertainty for those plug-in estimates. The interval
$[y_-,y_+]$ must be a genuine almost-sure support bound for the deterministic
remainder bound \eqref{eq:eta_m_sufficient}; if it is replaced by the observed
sample range, the result should be interpreted as a diagnostic rather than a
guaranteed population bound.
\end{remark}

% ======================================================================
\section{Proof}
\label{sec:proof}
% ======================================================================

\begin{proof}[Proof of Proposition~\ref{prop:abstract_bound}]
Throughout, $\|W\|_2=(\E[W^2])^{1/2}$. Write
\begin{equation}
  v=V-\E V,
  \qquad
  e=E-\E E,
  \qquad
  z=Z-\E Z=v+e.
\end{equation}
Then $\|v\|_2=\sigma_V$ and $\|e\|_2=\eta\sigma_V$.

For a nonempty collection $\mathscr C\subseteq\mathcal U_d$, define
\begin{equation}
  F_{\mathscr C}=\Pi_{\mathscr C}v,
  \qquad
  R_{\mathscr C}=\Pi_{\mathscr C}e.
\end{equation}
Because $\Pi_{\mathscr C}$ is an orthogonal projection, it is linear and
nonexpansive. Therefore
\begin{equation}
  \Pi_{\mathscr C}z=F_{\mathscr C}+R_{\mathscr C},
  \qquad
  \|R_{\mathscr C}\|_2\leq\|e\|_2=\eta\sigma_V.
  \label{eq:R_projection_bound}
\end{equation}
Let $p_{\mathscr C}=\mathsf{PE}_{\mathscr C}(V)$. Then
\begin{equation}
  \|F_{\mathscr C}\|_2^2=p_{\mathscr C}\sigma_V^2,
  \qquad
  \|F_{\mathscr C}\|_2=\sigma_V\sqrt{p_{\mathscr C}}.
  \label{eq:F_norm_general}
\end{equation}

First control the numerator projection. Expanding the squared norm gives
\begin{align}
  \|\Pi_{\mathscr C}z\|_2^2
  &=\|F_{\mathscr C}+R_{\mathscr C}\|_2^2
  \notag\\
  &=\|F_{\mathscr C}\|_2^2
    +2\langle F_{\mathscr C},R_{\mathscr C}\rangle
    +\|R_{\mathscr C}\|_2^2.
\end{align}
Hence, with
\begin{equation}
  \epsilon_N^{\mathscr C}
  =2\langle F_{\mathscr C},R_{\mathscr C}\rangle
   +\|R_{\mathscr C}\|_2^2,
\end{equation}
we have
\begin{equation}
  \|\Pi_{\mathscr C}z\|_2^2
  =p_{\mathscr C}\sigma_V^2+
    \epsilon_N^{\mathscr C}.
\end{equation}
Cauchy--Schwarz and \eqref{eq:R_projection_bound}--\eqref{eq:F_norm_general}
give
\begin{equation}
  |\epsilon_N^{\mathscr C}|
  \leq
  2\sigma_V\sqrt{p_{\mathscr C}}\,\eta\sigma_V
  +\eta^2\sigma_V^2
  =\sigma_V^2(2\eta\sqrt{p_{\mathscr C}}+\eta^2).
  \label{eq:epsN_general_bound}
\end{equation}

Next control the denominator. Since $z=v+e$,
\begin{equation}
  \Var[Z]=\|z\|_2^2=\|v+e\|_2^2.
\end{equation}
The reverse triangle inequality gives
\begin{equation}
  \|v+e\|_2\geq\bigl|\|v\|_2-\|e\|_2\bigr|
  =\sigma_V(1-\eta),
\end{equation}
because $\eta<1$. Therefore
\begin{equation}
  \Var[Z]\geq\sigma_V^2(1-\eta)^2>0.
  \label{eq:VarZ_general_lower}
\end{equation}
For the ratio identity, also write
\begin{equation}
  \Var[Z]=\sigma_V^2+
  \epsilon_D,
  \qquad
  \epsilon_D=2\langle v,e\rangle+\|e\|_2^2.
\end{equation}
Cauchy--Schwarz gives
\begin{equation}
  |\epsilon_D|
  \leq 2\sigma_V\eta\sigma_V+
  \eta^2\sigma_V^2
  =\sigma_V^2(2\eta+
  \eta^2).
  \label{eq:epsD_general_bound}
\end{equation}

Now subtract the two normalized projection energies over a common denominator:
\begin{align}
  \mathsf{PE}_{\mathscr C}(Z)-\mathsf{PE}_{\mathscr C}(V)
  &=\frac{\|\Pi_{\mathscr C}z\|_2^2}{\Var[Z]}
    -\frac{\|F_{\mathscr C}\|_2^2}{\sigma_V^2}
  \notag\\
  &=\frac{(p_{\mathscr C}\sigma_V^2+\epsilon_N^{\mathscr C})\sigma_V^2
          -p_{\mathscr C}\sigma_V^2(\sigma_V^2+
          \epsilon_D)}{\Var[Z]\sigma_V^2}
  \notag\\
  &=\frac{\epsilon_N^{\mathscr C}\sigma_V^2
          -p_{\mathscr C}\sigma_V^2\epsilon_D}{\Var[Z]\sigma_V^2}.
  \label{eq:general_ratio_identity}
\end{align}
Taking absolute values and applying \eqref{eq:epsN_general_bound},
\eqref{eq:epsD_general_bound}, and \eqref{eq:VarZ_general_lower},
\begin{align}
  \left|\mathsf{PE}_{\mathscr C}(Z)-\mathsf{PE}_{\mathscr C}(V)\right|
  &\leq
  \frac{|\epsilon_N^{\mathscr C}|+p_{\mathscr C}|\epsilon_D|}{\Var[Z]}
  \notag\\
  &\leq
  \frac{
    \sigma_V^2(2\eta\sqrt{p_{\mathscr C}}+\eta^2)
    +p_{\mathscr C}\sigma_V^2(2\eta+\eta^2)
  }{\sigma_V^2(1-\eta)^2}
  \notag\\
  &=
  \frac{
    2\eta\sqrt{p_{\mathscr C}}(1+\sqrt{p_{\mathscr C}})
    +\eta^2(1+p_{\mathscr C})
  }{(1-\eta)^2}.
\end{align}
This proves the proposition.
\end{proof}

\begin{proof}[Proof of Corollaries~\ref{cor:taylor_bound}--\ref{cor:special_cases}]
For Corollary~\ref{cor:taylor_bound}, Taylor's theorem gives
\begin{equation}
  \phi(Y)=\phi(\mu_Y)+V_m+R_m.
\end{equation}
Adding or subtracting the constant $\phi(\mu_Y)$ does not change any Sobol
index, so Proposition~\ref{prop:abstract_bound} applies with reference $V=V_m$
and perturbation $E=R_m$. The quantity $\eta_m$ is exactly the abstract
perturbation ratio $\eta$ from \eqref{eq:eta_def}. Taylor's theorem with
integral remainder gives
\begin{equation}
  R_m
  =\int_{\mu_Y}^{Y}\frac{(Y-t)^m}{m!}\phi^{(m+1)}(t)\,dt,
\end{equation}
and bounding $|\phi^{(m+1)}(t)|\leq\rho_{m+1}$ inside the integral and
evaluating yields the pointwise bound
\begin{equation}
  |R_m|
  \leq
  \frac{\rho_{m+1}}{(m+1)!}|Y-\mu_Y|^{m+1}.
  \label{eq:Rm_pointwise}
\end{equation}
(The integral form is preferred over the Lagrange form because the integration
variable $t$ runs over a deterministic interval once $Y$ is fixed; the Lagrange
form introduces an unknown evaluation point $\xi$ that depends on the
realization of $Y$ and cannot be tracked separately.) Hence
\begin{equation}
  \|R_m-\E R_m\|_2
  \leq \|R_m\|_2
  \leq
  \frac{\rho_{m+1}}{(m+1)!}
  \left\{\E|Y-\mu_Y|^{2m+2}\right\}^{1/2}.
\end{equation}
Dividing by $\sigma_m$ gives \eqref{eq:eta_m_sufficient}.

For Corollary~\ref{cor:local_affine}, take $m=1$. Then
$V_1=c(Y-\mu_Y)$, so affine invariance gives
$\mathsf{PE}_{\mathscr C}(V_1)=\mathsf{PE}_{\mathscr C}(Y)$ and
$\sigma_1=|c|\sigma_Y$. The Taylor remainder satisfies
\begin{equation}
  \|R_1-\E R_1\|_2
  \leq \|R_1\|_2
  \leq \frac12\rho_2\mu_4^{1/2},
\end{equation}
so
\begin{equation}
  \eta_1
  \leq
  \frac{\frac12\rho_2\mu_4^{1/2}}{|c|\sigma_Y}
  =\frac12\frac{\rho_2\sigma_Y}{|c|}\frac{\mu_4^{1/2}}{\sigma_Y^2}
  =\frac{\lambda}{2}.
\end{equation}
The bound in \eqref{eq:abstract_bound} is monotone increasing in $\eta$ on
$0\leq\eta<1$: the numerator is increasing and nonnegative, while the
denominator $(1-\eta)^2$ is decreasing and positive. Substituting the upper
bound $\eta_1\leq\lambda/2$ into \eqref{eq:abstract_bound} therefore gives
\eqref{eq:local_affine_bound}. Corollary~\ref{cor:special_cases} follows by
choosing the listed ANOVA projection collections.
\end{proof}

% ======================================================================
\section{Leading-order perturbation diagnostic}
\label{sec:perturbation}
% ======================================================================

The bound above is intentionally conservative: Cauchy--Schwarz is tight only
when the relevant projection of the reference output and the corresponding
projection of the residual are perfectly collinear. A first-order diagnostic
shows when the leading residual contribution cancels.

Let $V$ be any nonconstant square-integrable reference output and let
$G\in L^2$ be a candidate perturbation. Define
\begin{equation}
  Z_\varepsilon=V+\varepsilon G,
  \qquad
  v=V-\E V,
  \qquad
  g=G-\E G,
  \qquad
  \sigma_V^2=\Var[V].
\end{equation}
Let
\begin{equation}
  F_{\mathscr C}=\Pi_{\mathscr C}v,
  \qquad
  G_{\mathscr C}=\Pi_{\mathscr C}g,
  \qquad
  p_{\mathscr C}=\mathsf{PE}_{\mathscr C}(V).
\end{equation}
The numerator corresponding to $\mathsf{PE}_{\mathscr C}(Z_\varepsilon)$ is
$\|F_{\mathscr C}+\varepsilon G_{\mathscr C}\|_2^2$, while
\begin{equation}
  \Var[Z_\varepsilon]
  =\sigma_V^2+2\varepsilon\Cov[V,G]+\varepsilon^2\Var[G].
\end{equation}
Differentiating the ratio at $\varepsilon=0$ gives
\begin{equation}
  \left.\frac{d}{d\varepsilon}\mathsf{PE}_{\mathscr C}(Z_\varepsilon)\right|_{\varepsilon=0}
  =\frac{2}{\sigma_V^2}
  \left\{\Cov[F_{\mathscr C},G_{\mathscr C}]
         -p_{\mathscr C}\Cov[V,G]\right\}.
  \label{eq:directional_derivative_general}
\end{equation}
This formula is often more informative than the worst-case bound. Linear
cancellation occurs precisely when the perturbation covariance inside the chosen
ANOVA projection matches the denominator perturbation in the ratio
$p_{\mathscr C}$.

In the local-affine case $V=c(Y-\mu_Y)$ and $G=R_1$, this reduces to the earlier
curvature diagnostic. If $G(Y)=\frac12\tau(Y-\mu_Y)^2$ is the leading quadratic
Taylor perturbation, then for the first-order index of input $i$,
\begin{equation}
  \left.\frac{d}{d\varepsilon}S_i^{Z_\varepsilon}\right|_{\varepsilon=0}
  =\frac{\tau}{c\sigma_Y^2}
  \left\{\Cov[f_{\{i\}},\E[U^2\mid X_i]]-S_i^Y\E[U^3]\right\},
  \label{eq:quadratic_derivative}
\end{equation}
where $U=Y-\mu_Y$ and $f_{\{i\}}=\E[Y\mid X_i]-\mu_Y$. For an additive model
$Y=\mu_Y+\sum_j h_j(X_j)$ with independent centered components, this coefficient
is proportional to
\begin{equation}
  \E[h_i^3]-S_i^Y\sum_j\E[h_j^3].
  \label{eq:third_moment_contrast}
\end{equation}
Thus additivity alone does not force the leading perturbation to vanish;
symmetry or this precise third-moment contrast does.

For the \emph{critical-point regime}, where the reference is quadratic,
$V=\frac{b}{2}(Y-\mu_Y)^2$ with $b=\phi''(\mu_Y)\neq0$, the corresponding
diagnostic uses $G=R_2$ and the general formula \eqref{eq:directional_derivative_general}
with $\mathscr C=\{i\}$:
\begin{equation}
  \left.\frac{d}{d\varepsilon}S_i^{Z_\varepsilon}\right|_{\varepsilon=0}
  =\frac{2}{\sigma_{V}^2}
  \left\{\Cov\!\left[\Pi_{\{i\}}V^\circ,\,\Pi_{\{i\}}G^\circ\right]
         -p_{\{i\}}\Cov[V,G]\right\},
  \label{eq:quadratic_ref_derivative}
\end{equation}
where $V^\circ=V-\E V$ and $G^\circ=G-\E G$. The denominator
$\sigma_V^2=\Var[V]$ is the variance of the quadratic reference, not a squared
slope. Under the required condition $\Var[V]>0$, this formula does not involve
division by $\phi'(\mu_Y)$ and applies without any nonzero-slope requirement. When $G=R_2$ is the order-3 Taylor residual, the leading
coefficient is proportional to the covariance between the first-order ANOVA projection of the
quadratic reference and the first-order projection of the cubic remainder, minus
the index value times the corresponding denominator perturbation.

% ======================================================================
\section{Canonical quadratic example}
\label{sec:numerical}
% ======================================================================

Let
\begin{equation}
  Y=X_1+X_2,
  \qquad
  X_1,X_2\overset{\mathrm{iid}}{\sim}\mathrm{Uniform}(-1,1),
  \qquad
  Z=\phi(Y)=Y^2.
\end{equation}
Then $\mu_Y=0$ and $\phi'(\mu_Y)=0$, so any theorem that requires a nonzero
slope at the mean excludes the canonical mechanism. The unified theorem does not
exclude it. Choose $m=2$. Since $\phi(y)=y^2$ is exactly quadratic,
\begin{equation}
  V_2=Y^2,
  \qquad
  R_2=0,
  \qquad
  \eta_2=0.
\end{equation}
The bound therefore says that the indices of $Z$ equal the indices of the
quadratic reference $V_2$ exactly.

The ANOVA structure of that reference displays the mechanism. Expanding,
\begin{equation}
  Z=(X_1+X_2)^2=X_1^2+2X_1X_2+X_2^2.
\end{equation}
Because $\E[X_i]=0$, $\E[X_i^2]=1/3$, and $\E[X_i^4]=1/5$, the ANOVA components
are
\begin{equation}
  z_1(X_1)=X_1^2-\tfrac13,
  \qquad
  z_2(X_2)=X_2^2-\tfrac13,
  \qquad
  z_{12}(X_1,X_2)=2X_1X_2.
\end{equation}
Their variance contributions are
\begin{equation}
  D_1^Z=D_2^Z=\Var(X_i^2)=\tfrac15-\tfrac19=\tfrac4{45},
  \qquad
  D_{12}^Z=\Var(2X_1X_2)=4\E[X_1^2]\E[X_2^2]=\tfrac49.
\end{equation}
Therefore
\begin{equation}
  \Var[Z]=\tfrac4{45}+\tfrac4{45}+\tfrac49=\tfrac{28}{45},
\end{equation}
and
\begin{equation}
  S_1^Z=S_2^Z=\frac{1}{7},
  \qquad
  S_{\{1,2\}}^Z=\frac{5}{7},
  \qquad
  (S_1^T)^Z=(S_2^T)^Z=\frac{6}{7}.
  \label{eq:quadratic_indices}
\end{equation}
The original output $Y=X_1+X_2$ is additive, with $S_1^Y=S_2^Y=1/2$ and
$S_{\{1,2\}}^Y=0$. Squaring does not merely perturb the original additive ANOVA
structure; it creates a dominant interaction term. This is why the bound must be
formulated relative to the Taylor reference $V_m$ rather than always relative to
$Y$.

For contrast, if $\phi'(\mu_Y)\neq0$ and $m=1$, the reference is affine in $Y$
and Corollary~\ref{cor:local_affine} recovers a direct bound on shifts relative
to the Sobol indices of $Y$. If $\phi$ is a low-degree polynomial and $m$ is
chosen at least as large as its degree, the residual is zero and the theorem is
exact; the work then lies in understanding the Sobol decomposition of the
polynomial reference.

% ======================================================================
\appendix

\section{Pedagogical proof: every step explained}
\label{app:pedagogical}
% ======================================================================

This appendix re-derives Proposition~\ref{prop:abstract_bound} and its
Taylor-reference corollaries from first principles. The goal is not brevity. The
goal is to make explicit what each symbol means, why projections appear, why the
denominator needs a stability condition, and how the Taylor approximation enters
only after the abstract projection argument is finished.

\subsection*{Orientation before the algebra}

Before beginning the line-by-line proof, it is useful to say what the proof is
trying to accomplish and why it is organized in this somewhat abstract way. The
problem begins with a practical question: after replacing an output $Y$ by a
transformed output $Z=\phi(Y)$, how much can the usual variance-based Sobol
indices move? A tempting first answer is to linearize $\phi$ and compare
$\phi(Y)$ directly with $Y$. That answer is incomplete. It works when the
linear Taylor term dominates, but it misses the most important failure mode of
commutativity: a nonlinear transform can change the ANOVA structure of the
output. The centered squaring example is the canonical case. If
$Y=X_1+X_2$ and $\phi(y)=y^2$, then
\begin{equation}
  \phi(Y)=(X_1+X_2)^2=X_1^2+2X_1X_2+X_2^2,
\end{equation}
so the transformed output contains a genuine interaction even though the
original output is additive. A proof that only compares $\phi(Y)$ with $Y$
would either exclude this example or obscure the mechanism that makes it
important.

The proof therefore separates two tasks. The first task is purely geometric: if
two centered outputs are close in $L^2$, then their projected variance shares are
close, provided the denominator variance of the perturbed output does not nearly
vanish. This part has nothing to do with Taylor expansions, smoothness, or the
value of $\phi'(\mu_Y)$. It is a Hilbert-space perturbation argument for outputs
of the form
\begin{equation}
  Z=V+E,
\end{equation}
where $V$ is a reference output and $E$ is a residual. The second task is to use
Taylor's theorem to choose a scientifically meaningful reference $V$. For a
smooth transformation, the reference is the Taylor polynomial in the centered
output
\begin{equation}
  V_m=\sum_{k=1}^{m}\frac{\phi^{(k)}(\mu_Y)}{k!}(Y-\mu_Y)^k,
\end{equation}
and the residual is the Taylor remainder. This division of labor is the main
structural point of the memo. The abstract proposition says what any small
output perturbation can do to Sobol variance shares; the Taylor corollaries say
which reference output should be used for a transformed scalar response.

This organization also explains why the theorem is not a second commutativity
theorem. The companion commutativity result asks when the Sobol indices of
$\phi(Y)$ must equal the Sobol indices of $Y$ for every admissible model. The
answer is: only for affine transformations. The present result asks a different,
quantitative question. It asks when the Sobol indices of $\phi(Y)$ are close to
the Sobol indices of the dominant Taylor reference $V_m$. In the local-affine
case, $V_m$ may be an affine rescaling of $Y$, so the comparison is directly to
$Y$. In the critical-point case, the first nonzero Taylor term may be quadratic
or higher-order, so the comparison is to that polynomial reference instead. This
is not a weakness of the proof; it is what lets the theorem include the
interaction-generating examples rather than excluding them.

The main technical challenge is that a Sobol index is a ratio. Its numerator is
a variance contribution assigned to one ANOVA component, a closed collection of
components, or a total-effect collection. Its denominator is the total variance
of the output. A small residual can perturb both pieces. The numerator is
manageable because Sobol numerators are squared lengths of orthogonal
projections. Projection is linear, and orthogonal projections are contractions:
projecting a residual cannot make its $L^2$ norm larger. Thus, when
$Z=V+E$, the projected numerator for $Z$ is the projected numerator for $V$ plus
a cross term and a residual-squared term. Cauchy--Schwarz controls the cross
term, and contraction controls the residual term. This is where the
$2\eta\sqrt{p}$ and $\eta^2$ pieces of the bound come from.

The denominator is more delicate because ratios can become unstable when the
denominator is small. Even if $E$ is small in norm relative to $V$, the centered
residual could point partly against the centered reference and reduce the total
variance of $Z$. The proof therefore imposes the dimensionless stability
condition
\begin{equation}
  \eta=\frac{\|E-\E E\|_2}{\operatorname{sd}(V)}<1.
\end{equation}
The reverse triangle inequality then gives
\begin{equation}
  \operatorname{sd}(Z)\geq \operatorname{sd}(V)(1-\eta),
\end{equation}
which is enough to keep the denominator away from zero. This step is deliberately
conservative. It does not try to exploit the sign of a covariance between the
reference and residual. The payoff is a clean bound that works uniformly for any
ANOVA projection collection.

A useful way to read the proof is to track three quantities. First,
\begin{equation}
  p_{\mathscr C}=\mathsf{PE}_{\mathscr C}(V)
\end{equation}
is the reference share of variance assigned to the ANOVA collection
$\mathscr C$. Second,
\begin{equation}
  \eta=\frac{\|E-\E E\|_2}{\sigma_V}
\end{equation}
measures residual size relative to the standard deviation of the reference. This
is the perturbation parameter. Third, $(1-\eta)^2$ is the denominator stability
factor. The final bound is just the numerator perturbation divided by the stable
lower bound for the denominator, plus the correction needed to compare two
ratios rather than two raw projected variances.

The proof is also intentionally stated for a generic projection collection
$\mathscr C$. This is a notational convenience, not a new sensitivity index. If
$\mathscr C=\{u\}$, then the projection energy is the Sobol ANOVA component
index $S_u$. If $\mathscr C$ contains all nonempty subsets of a group $A$, it is
the closed index $(S_A^{\mathrm{cl}})$. If $\mathscr C$ contains all subsets
that intersect $A$, it is the total-effect index $(S_A^T)$. The same algebra
works for all three because all three are squared projection norms divided by
total variance. This avoids summing many componentwise bounds and avoids adding
artificial factors depending on the number of subsets.

The proof has limitations that should be understood before reading the details.
It is a population result, not a sampling theorem. Estimating
$\eta$, $p_{\mathscr C}$, or the Taylor-reference Sobol indices requires its own
Monte Carlo or statistical uncertainty analysis. The Taylor corollary also
requires the reference $V_m$ to have positive variance. A nonzero Taylor
coefficient is not enough by itself: on a degenerate or finite support, the
polynomial reference could still be constant. Finally, the residual must be
small relative to the chosen reference for the bound to be useful. If the Taylor
remainder is large, the theorem correctly reports that the chosen reference does
not explain the transformed output well.

With that overview in mind, the proof proceeds in eight steps. Step~0 sets up
the $L^2$ language and explains why variance is squared norm after centering.
Step~1 rewrites Sobol component, closed, and total-effect indices as normalized
projection energies. Step~2 states the abstract perturbation decomposition
$Z=V+E$ and defines the dimensionless residual ratio. Step~3 controls the
projected numerator. Step~4 controls the denominator. Step~5 combines the two
controls into the projection-energy bound. Step~6 translates the generic
projection-energy result back into standard Sobol notation. Step~7 inserts
Taylor's theorem and obtains an explicit residual bound. Step~8 interprets the
special cases: the local-affine case, the critical-point quadratic case, and the
exact-polynomial case. The detailed algebra below is therefore not a collection
of unrelated inequalities; each inequality serves one of these three purposes:
control projected numerator, stabilize denominator, or connect the abstract
reference-residual argument to a Taylor approximation of $\phi(Y)$.

\subsection*{Step 0. The probability and Hilbert-space setup}

The inputs are
\begin{equation}
  \bX=(X_1,\ldots,X_d),
\end{equation}
and the components are independent. Independence is what makes the
Hoeffding--Sobol ANOVA decomposition orthogonal. The output of a deterministic
model is a scalar random variable because the random input vector is passed
through a deterministic function:
\begin{equation}
  W=f(\bX).
\end{equation}
The phrase $W\in L^2$ means that $W$ has a finite second moment,
\begin{equation}
  \E[W^2]<\infty.
\end{equation}
For two square-integrable random variables $A$ and $B$, define
\begin{equation}
  \langle A,B\rangle=\E[AB],
  \qquad
  \|A\|_2=(\E[A^2])^{1/2}.
\end{equation}
This is the usual inner-product geometry of $L^2$. If $W^\circ=W-\E W$, then
\begin{equation}
  \|W^\circ\|_2^2
  =\E[(W-\E W)^2]
  =\Var[W].
\end{equation}
Thus, after centering, squared $L^2$ norm is variance and $L^2$ norm is standard
deviation.

This centering step is essential. Sobol indices allocate variance, not raw second
moment. Adding a constant to an output changes $W$ but does not change
$W-\E W$, and therefore does not change any Sobol numerator or denominator.
This is why constants such as $\phi(\mu_Y)$ disappear from the proof.

\subsection*{Step 1. Sobol components as orthogonal pieces of a centered output}

For independent inputs, any square-integrable scalar output $W$ has a unique
Hoeffding--Sobol decomposition
\begin{equation}
  W=\E W+\sum_{u\in\mathcal U_d}w_u(\bX_u),
  \qquad
  \mathcal U_d={u\subseteq{1,\ldots,d}:u\neq\emptyset}.
\end{equation}
Here $u$ is an input subset. For example, $u={1}$ is a main-effect component,
while $u={1,2}$ is a two-way interaction component. The notation
$\bX_u=(X_j)_{j\in u}$ means the subvector of inputs indexed by $u$.

Each component $w_u$ has mean zero and represents the part of the centered output
associated with the subset $u$. The components are mutually orthogonal:
\begin{equation}
  \E[w_u(\bX_u)w_v(\bX_v)]=0,
  \qquad u\neq v.
\end{equation}
Orthogonality gives the variance decomposition
\begin{align}
  \Var[W]
  &=\left\|\sum_{u\in\mathcal U_d}w_u\right\|_2^2 \\
  &=\sum_{u\in\mathcal U_d}\|w_u\|_2^2.
\end{align}
The variance contribution of component $u$ is
\begin{equation}
  D_u^W=\|w_u\|_2^2,
\end{equation}
and the Sobol ANOVA component index is
\begin{equation}
  S_u^W=\frac{D_u^W}{\Var[W]}
       =\frac{\|w_u\|_2^2}{\Var[W]}.
\end{equation}

The proof also needs closed and total-effect indices. Rather than proving the
same inequality three times, collect several ANOVA subspaces at once. For a
nonempty collection $\mathscr C\subseteq\mathcal U_d$, define
\begin{equation}
  \Pi_{\mathscr C}(W-\E W)=\sum_{u\in\mathscr C}w_u(\bX_u).
\end{equation}
The symbol $\Pi_{\mathscr C}$ means orthogonal projection onto the direct sum of
the ANOVA subspaces indexed by $\mathscr C$. Because the selected components are
orthogonal,
\begin{align}
  \|\Pi_{\mathscr C}(W-\E W)\|_2^2
  &=\left\|\sum_{u\in\mathscr C}w_u\right\|_2^2 \\
  &=\sum_{u\in\mathscr C}\|w_u\|_2^2.
\end{align}
Therefore
\begin{equation}
  \mathsf{PE}_{\mathscr C}(W)
  =\frac{\|\Pi_{\mathscr C}(W-\E W)\|_2^2}{\Var[W]}
  =\sum_{u\in\mathscr C}S_u^W.
\end{equation}
This notation is only normalized projection energy. It is not a new sensitivity
measure.

The standard cases are recovered by choosing $\mathscr C$:
\begin{itemize}
  \item component index: choose $\mathscr C={u}$, giving
    $\mathsf{PE}_{{u}}(W)=S_u^W$;
  \item closed index for an input group $A$: choose
    $\mathscr C={u:\emptyset\neq u\subseteq A}$, giving
    $(S_A^{\mathrm{cl}})^W$;
  \item total-effect index for $A$: choose
    $\mathscr C={u:u\cap A\neq\emptyset}$, giving $(S_A^T)^W$.
\end{itemize}
This is why a single projection theorem can cover component, closed, and
total-effect Sobol indices without summing many separately bounded component
errors.

\subsection*{Step 2. Split the transformed output into a reference and a residual}

The abstract theorem starts with two square-integrable outputs $V$ and $E$ and
writes
\begin{equation}
  Z=V+E.
\end{equation}
Here $V$ is the reference output: the part whose Sobol structure we understand or
want to compare against. The term $E$ is the residual perturbation.

Because Sobol indices use centered outputs, define
\begin{equation}
  v=V-\E V,
  \qquad
  e=E-\E E,
  \qquad
  z=Z-\E Z.
\end{equation}
Then
\begin{align}
  z
  &=Z-\E Z \\
  &=(V+E)-\E[V+E] \\
  &=(V-\E V)+(E-\E E) \\
  &=v+e.
\end{align}
Let
\begin{equation}
  \sigma_V^2=\Var[V]=\|v\|_2^2>0,
  \qquad
  \eta=\frac{\|e\|_2}{\sigma_V}.
\end{equation}
The number $\eta$ is dimensionless. It measures the centered residual size in
units of the reference standard deviation. The assumption $\eta<1$ says that the
residual is smaller than the reference in $L^2$ norm. It does not say the
residual is pointwise smaller for every input realization; it is a variance-scale
condition.

\subsection*{Step 3. Why projections cannot amplify the residual}

Fix a collection $\mathscr C$ of ANOVA subsets and define
\begin{equation}
  F_{\mathscr C}=\Pi_{\mathscr C}v,
  \qquad
  R_{\mathscr C}=\Pi_{\mathscr C}e.
\end{equation}
The letter $F$ denotes the selected ANOVA part of the reference, while $R$
denotes the selected ANOVA part of the residual. Projection linearity gives
\begin{equation}
  \Pi_{\mathscr C}z
  =\Pi_{\mathscr C}(v+e)
  =F_{\mathscr C}+R_{\mathscr C}.
\end{equation}

An orthogonal projection is nonexpansive. To see this, write any element $a$ as
\begin{equation}
  a=\Pi_{\mathscr C}a+(I-\Pi_{\mathscr C})a,
\end{equation}
where the two pieces are orthogonal. Pythagoras gives
\begin{equation}
  \|a\|_2^2
  =\|\Pi_{\mathscr C}a\|_2^2+\|(I-\Pi_{\mathscr C})a\|_2^2
  \geq \|\Pi_{\mathscr C}a\|_2^2.
\end{equation}
Thus
\begin{equation}
  \|\Pi_{\mathscr C}a\|_2\leq\|a\|_2.
\end{equation}
Applying this to $a=e$ gives
\begin{equation}
  \|R_{\mathscr C}\|_2
  =\|\Pi_{\mathscr C}e\|_2
  \leq\|e\|_2
  =\eta\sigma_V.
\end{equation}
The residual may have complicated ANOVA structure, but no projection of it can
have more $L^2$ energy than the full residual.

The reference projection energy is
\begin{equation}
  p_{\mathscr C}=\mathsf{PE}_{\mathscr C}(V)
  =\frac{\|F_{\mathscr C}\|_2^2}{\sigma_V^2}.
\end{equation}
Therefore
\begin{equation}
  \|F_{\mathscr C}\|_2=\sigma_V\sqrt{p_{\mathscr C}}.
\end{equation}
This identity is where the reference Sobol index enters the bound.

\subsection*{Step 4. Expand and bound the projected numerator}

The numerator of $\mathsf{PE}_{\mathscr C}(Z)$ is the selected projected variance
of $Z$:
\begin{equation}
  \|\Pi_{\mathscr C}z\|_2^2.
\end{equation}
Using the decomposition from Step 3,
\begin{align}
  \|\Pi_{\mathscr C}z\|_2^2
  &=\|F_{\mathscr C}+R_{\mathscr C}\|_2^2 \\
  &=\|F_{\mathscr C}\|_2^2
    +2\langle F_{\mathscr C},R_{\mathscr C}\rangle
    +\|R_{\mathscr C}\|_2^2.
\end{align}
Define the numerator perturbation
\begin{equation}
  \epsilon_N^{\mathscr C}
  =2\langle F_{\mathscr C},R_{\mathscr C}\rangle
   +\|R_{\mathscr C}\|_2^2.
\end{equation}
Then
\begin{equation}
  \|\Pi_{\mathscr C}z\|_2^2
  =p_{\mathscr C}\sigma_V^2+
  \epsilon_N^{\mathscr C}.
\end{equation}

Now bound $\epsilon_N^{\mathscr C}$. The Cauchy--Schwarz inequality says
\begin{equation}
  |\langle A,B\rangle|\leq\|A\|_2\|B\|_2.
\end{equation}
Hence
\begin{align}
  |\langle F_{\mathscr C},R_{\mathscr C}\rangle|
  &\leq \|F_{\mathscr C}\|_2\|R_{\mathscr C}\|_2 \\
  &\leq (\sigma_V\sqrt{p_{\mathscr C}})(\eta\sigma_V) \\
  &=\eta\sigma_V^2\sqrt{p_{\mathscr C}}.
\end{align}
Also,
\begin{equation}
  \|R_{\mathscr C}\|_2^2\leq\eta^2\sigma_V^2.
\end{equation}
Therefore
\begin{align}
  |\epsilon_N^{\mathscr C}|
  &\leq 2\eta\sigma_V^2\sqrt{p_{\mathscr C}}
       +\eta^2\sigma_V^2 \\
  &=\sigma_V^2(2\eta\sqrt{p_{\mathscr C}}+
       \eta^2).
  \label{eq:ped_general_epsN}
\end{align}
The two terms have different interpretations. The term linear in $\eta$ is the
cross term between the reference and residual projections. The term quadratic in
$\eta$ is the projected variance of the residual itself.

\subsection*{Step 5. Bound the denominator away from zero}

The denominator of $\mathsf{PE}_{\mathscr C}(Z)$ is
\begin{equation}
  \Var[Z]=\|z\|_2^2=\|v+e\|_2^2.
\end{equation}
We need a lower bound because a small numerator difference can produce a large
ratio difference if the transformed-output variance is nearly zero.

The reverse triangle inequality follows from the ordinary triangle inequality:
\begin{equation}
  \|v\|_2=\|(v+e)-e\|_2\leq\|v+e\|_2+\|e\|_2,
\end{equation}
so
\begin{equation}
  \|v+e\|_2\geq\|v\|_2-\|e\|_2.
\end{equation}
Since $\|v\|_2=\sigma_V$ and $\|e\|_2=\eta\sigma_V$,
\begin{equation}
  \|v+e\|_2\geq\sigma_V(1-\eta).
\end{equation}
When $\eta<1$, the right side is positive. Squaring gives
\begin{equation}
  \Var[Z]\geq\sigma_V^2(1-\eta)^2>0.
  \label{eq:ped_general_var_lower}
\end{equation}
This proves that the transformed output is nonconstant, so its Sobol indices are
well-defined.

We also need to compare $\Var[Z]$ to $\Var[V]$. Expanding gives
\begin{align}
  \Var[Z]
  &=\|v+e\|_2^2 \\
  &=\|v\|_2^2+2\langle v,e\rangle+\|e\|_2^2 \\
  &=\sigma_V^2+\epsilon_D,
\end{align}
where
\begin{equation}
  \epsilon_D=2\langle v,e\rangle+\|e\|_2^2.
\end{equation}
Here $\epsilon_D$ is the denominator perturbation. Cauchy--Schwarz gives
\begin{align}
  |\epsilon_D|
  &\leq 2\|v\|_2\|e\|_2+
       \|e\|_2^2 \\
  &=2\sigma_V(\eta\sigma_V)+\eta^2\sigma_V^2 \\
  &=\sigma_V^2(2\eta+
       \eta^2).
  \label{eq:ped_general_epsD}
\end{align}
This is the denominator analogue of the numerator bound.

\subsection*{Step 6. Subtract the two Sobol ratios}

Now write the two normalized projection energies explicitly:
\begin{equation}
  \mathsf{PE}_{\mathscr C}(Z)
  =\frac{p_{\mathscr C}\sigma_V^2+
  \epsilon_N^{\mathscr C}}{\Var[Z]},
  \qquad
  \mathsf{PE}_{\mathscr C}(V)
  =\frac{p_{\mathscr C}\sigma_V^2}{\sigma_V^2}
  =p_{\mathscr C}.
\end{equation}
Using $\Var[Z]=\sigma_V^2+
\epsilon_D$, put the difference over a common denominator:
\begin{align}
  \mathsf{PE}_{\mathscr C}(Z)-\mathsf{PE}_{\mathscr C}(V)
  &=\frac{p_{\mathscr C}\sigma_V^2+
      \epsilon_N^{\mathscr C}}{\Var[Z]}
    -\frac{p_{\mathscr C}\sigma_V^2}{\sigma_V^2} \\
  &=\frac{(p_{\mathscr C}\sigma_V^2+
      \epsilon_N^{\mathscr C})\sigma_V^2
      -p_{\mathscr C}\sigma_V^2(\sigma_V^2+
      \epsilon_D)}{\Var[Z]\sigma_V^2} \\
  &=\frac{\cancel{p_{\mathscr C}\sigma_V^4}
      +\epsilon_N^{\mathscr C}\sigma_V^2
      -\cancel{p_{\mathscr C}\sigma_V^4}
      -p_{\mathscr C}\sigma_V^2\epsilon_D}{\Var[Z]\sigma_V^2} \\
  &=\frac{\epsilon_N^{\mathscr C}
      -p_{\mathscr C}\epsilon_D}{\Var[Z]}.
  \label{eq:ped_general_ratio}
\end{align}
The reference contribution cancels. This cancellation is the ratio-level reason
that the result is scale-free: multiplying $V$ and $E$ by a common constant would
not change either Sobol index.

Taking absolute values and applying the triangle inequality,
\begin{equation}
  \left|\mathsf{PE}_{\mathscr C}(Z)-\mathsf{PE}_{\mathscr C}(V)\right|
  \leq
  \frac{|\epsilon_N^{\mathscr C}|+p_{\mathscr C}|\epsilon_D|}{\Var[Z]}.
\end{equation}
Substitute the numerator bound \eqref{eq:ped_general_epsN}, the denominator
perturbation bound \eqref{eq:ped_general_epsD}, and the variance lower bound
\eqref{eq:ped_general_var_lower}:
\begin{align}
  \left|\mathsf{PE}_{\mathscr C}(Z)-\mathsf{PE}_{\mathscr C}(V)\right|
  &\leq
  \frac{
    \sigma_V^2(2\eta\sqrt{p_{\mathscr C}}+
    \eta^2)
    +p_{\mathscr C}\sigma_V^2(2\eta+
    \eta^2)
  }{\sigma_V^2(1-\eta)^2} \\
  &=
  \frac{
    2\eta\sqrt{p_{\mathscr C}}+
    \eta^2+2\eta p_{\mathscr C}+
    \eta^2p_{\mathscr C}
  }{(1-\eta)^2} \\
  &=
  \frac{
    2\eta\sqrt{p_{\mathscr C}}(1+\sqrt{p_{\mathscr C}})
    +\eta^2(1+p_{\mathscr C})
  }{(1-\eta)^2}.
\end{align}
This is Proposition~\ref{prop:abstract_bound}.

\subsection*{Step 7. Taylor expansion supplies a concrete reference}

The abstract theorem says nothing about derivatives. It only needs a reference
$V$ and residual $E$. For a smooth output transformation $Z=\phi(Y)$, Taylor
expansion provides a principled choice.

Let
\begin{equation}
  \mu_Y=\E[Y],
  \qquad
  U=Y-\mu_Y.
\end{equation}
Taylor expansion around $\mu_Y$ through order $m$ gives
\begin{equation}
  \phi(Y)=\phi(\mu_Y)+
  \sum_{k=1}^m\frac{\phi^{(k)}(\mu_Y)}{k!}U^k+R_m.
\end{equation}
The constant $\phi(\mu_Y)$ has no effect on Sobol indices. Therefore choose
\begin{equation}
  V_m=\sum_{k=1}^m\frac{\phi^{(k)}(\mu_Y)}{k!}U^k,
  \qquad
  E=R_m.
\end{equation}
The reference variance condition is
\begin{equation}
  \sigma_m^2=\Var[V_m]>0.
\end{equation}
This condition is real. It is not automatically guaranteed by a nonzero Taylor
coefficient, because a nonzero polynomial can still be constant on the particular
support of $Y$. For example, if $U$ takes only the two values $-1$ and $1$, then
$U^2=1$ almost surely and a purely quadratic reference has zero variance.

The Taylor residual ratio is
\begin{equation}
  \eta_m=\frac{\|R_m-\E R_m\|_2}{\sigma_m}.
\end{equation}
If $\eta_m<1$, the abstract theorem directly bounds Sobol indices of $\phi(Y)$
relative to Sobol indices of $V_m$.

Now derive the computable sufficient bound. Let
\begin{equation}
  \rho_{m+1}=\|\phi^{(m+1)}\|_{\infty,[y_-,y_+]}.
\end{equation}
The integral form of the Taylor remainder is
\begin{equation}
  R_m
  =\int_{\mu_Y}^{Y}
    \frac{(Y-t)^m}{m!}\phi^{(m+1)}(t)\,dt.
\end{equation}
Taking absolute values and using
$|\phi^{(m+1)}(t)|\leq\rho_{m+1}$ gives
\begin{align}
  |R_m|
  &\leq
  \frac{\rho_{m+1}}{m!}
  \left|\int_{\mu_Y}^{Y}(Y-t)^m\,dt\right| \\
  &=
  \frac{\rho_{m+1}}{m!}
  \frac{|Y-\mu_Y|^{m+1}}{m+1} \\
  &=
  \frac{\rho_{m+1}}{(m+1)!}|Y-\mu_Y|^{m+1}.
\end{align}
Squaring, taking expectations, and taking square roots gives
\begin{equation}
  \|R_m\|_2
  \leq
  \frac{\rho_{m+1}}{(m+1)!}
  \left\{\E|Y-\mu_Y|^{2m+2}\right\}^{1/2}.
\end{equation}
Finally,
\begin{equation}
  \|R_m-\E R_m\|_2^2
  =\Var[R_m]
  =\E[R_m^2]-(\E R_m)^2
  \leq \E[R_m^2]
  =\|R_m\|_2^2.
\end{equation}
Therefore
\begin{equation}
  \eta_m
  \leq
  \frac{\rho_{m+1}}{(m+1)!\,\sigma_m}
  \left\{\E|Y-\mu_Y|^{2m+2}\right\}^{1/2},
\end{equation}
which is the sufficient bound in Corollary~\ref{cor:taylor_bound}.

\subsection*{Step 8. How the theorem reads in the familiar regimes}

If $m=1$ and $c=\phi'(\mu_Y)\neq0$, then
\begin{equation}
  V_1=c(Y-\mu_Y).
\end{equation}
This is an affine rescaling of $Y$, after removing an irrelevant constant. Its
ANOVA components are $c$ times the ANOVA components of $Y$, so each component
variance and the total variance are multiplied by $c^2$. The ratios are
unchanged:
\begin{equation}
  \mathsf{PE}_{\mathscr C}(V_1)=\mathsf{PE}_{\mathscr C}(Y).
\end{equation}
Thus the general theorem becomes a direct bound on index shifts relative to the
original output $Y$.

If $\phi'(\mu_Y)=0$ but $\phi''(\mu_Y)\neq0$, the first useful reference is often
quadratic:
\begin{equation}
  V_2=\frac12\phi''(\mu_Y)(Y-\mu_Y)^2.
\end{equation}
The theorem applies when $\Var[V_2]>0$. It does not compare $\phi(Y)$ directly to
$Y$; it compares $\phi(Y)$ to the quadratic reference. That distinction is the
point. If $Y=X_1+X_2$ is additive and centered, then $Y^2$ contains
$2X_1X_2$, so the quadratic reference has an interaction component even though
$Y$ does not.

More generally, choose $m$ large enough that the residual $R_m$ is small relative
to the standard deviation of $V_m$. The theorem then says that the Sobol
structure of the transformed output is controlled by the Sobol structure of the
Taylor reference plus a quantified residual error. The proof is unified because
all regimes use the same abstract decomposition $Z=V+E$.

% ======================================================================
\section{Proof story: from question to finished argument}
\label{app:proof_story}
% ======================================================================

\subsection*{Starting point: the commutativity memo}

The companion memo says that affine transformations are exactly the continuous
pointwise transformations that commute with Sobol index computation for every
admissible scalar model. That result is qualitative. It tells us that a
non-affine transformation can move Sobol indices, but it does not tell us how
large the movement can be for a particular model and transform.

The first instinct is to linearize $\phi$ around $\mu_Y$ and compare $\phi(Y)$ to
$Y$. That works when $\phi'(\mu_Y)\neq0$, because the leading Taylor term is an
affine rescaling of $Y-\mu_Y$. But it fails as a general organizing principle:
if $\phi'(\mu_Y)=0$, the leading term may be quadratic or higher-order. The
canonical example $Y=X_1+X_2$ and $\phi(y)=y^2$ is exactly this case; the
quadratic term creates the interaction $2X_1X_2$.

The fix is to separate two ideas that were previously entangled. The projection
perturbation theorem should not know anything about Taylor expansions or about
whether the first derivative is zero. It should only say: if $Z=V+E$ and $E$ is
small relative to $V$, then Sobol indices of $Z$ are close to Sobol indices of
$V$. Taylor analysis is then used afterward to choose $V$.

\subsection*{Why the reference output is allowed to change}

Sobol indices are not intrinsic properties of an input--output mechanism alone;
they are variance shares on a chosen output scale. If the output is squared, the
variance being allocated is the variance of the squared output. There is no
mathematical reason that this variance allocation must resemble the variance
allocation of the unsquared output.

The reference $V_m$ is therefore not a trick. It is the dominant output-scale
object after transformation. When the dominant Taylor term is affine, $V_m$ has
the same Sobol indices as $Y$. When the dominant Taylor term is quadratic,
$V_m$ may have new interactions. The theorem follows the mathematics of the
transformed output rather than forcing every transform into a local-affine frame.

\subsection*{Why orthogonal projections are the right language}

A Sobol component numerator is the squared norm of an orthogonal projection of
the centered output onto an ANOVA subspace. A closed index or total-effect index
is the squared norm of a projection onto a direct sum of such subspaces. Once
this is recognized, all standard Sobol indices in the memo have the same form:
\begin{equation}
  \frac{\|\Pi_{\mathscr C}(W-\E W)\|_2^2}{\Var[W]}.
\end{equation}
This is why the theorem is stated for $\mathsf{PE}_{\mathscr C}$. The notation
is bookkeeping: it lets one proof handle component, closed, and total-effect
indices without repeating the same calculation and without summing many separate
componentwise bounds.

\subsection*{The numerator argument}

If $Z=V+E$, then after centering,
\begin{equation}
  Z-\E Z=(V-\E V)+(E-\E E)=v+e.
\end{equation}
Projection linearity gives
\begin{equation}
  \Pi_{\mathscr C}(Z-\E Z)=\Pi_{\mathscr C}v+\Pi_{\mathscr C}e.
\end{equation}
The difference between the projected numerator for $Z$ and the projected
numerator for $V$ is therefore just the cross term plus the squared perturbation
projection. Cauchy--Schwarz bounds the cross term, and projection contraction
bounds the perturbation projection by the full perturbation norm. This gives a
linear term in $\eta$ and a quadratic term in $\eta$.

\subsection*{The denominator argument}

The denominator is the total variance of $Z$, namely $\|v+e\|_2^2$. The residual
could partially cancel the reference, so a lower bound is needed. The reverse
triangle inequality gives the clean stability condition
\begin{equation}
  \|v+e\|_2\geq \|v\|_2-\|e\|_2=\sigma_V(1-\eta).
\end{equation}
Thus $\eta<1$ guarantees that the transformed output is nonconstant and that the
ratio denominator is bounded away from zero. This is conceptually simpler than
tracking signs of covariance terms, and it applies to every reference $V$, not
just an affine one.

\subsection*{The key cancellation}

The ratio subtraction is where the proof becomes a Sobol-index proof rather than
just an $L^2$ perturbation bound. Writing both projection energies over a common
denominator makes the reference contribution cancel. What remains is the
numerator perturbation minus the reference projection share times the denominator
perturbation. This cancellation is why the bound depends on the reference index
value $p_{\mathscr C}$ and on the relative residual size $\eta$, but not on the
absolute scale of $V$.

\subsection*{How Taylor enters}

Taylor expansion provides a principled way to choose $V$. Around $\mu_Y$,
\begin{equation}
  \phi(Y)=\phi(\mu_Y)+V_m+R_m,
  \qquad
  V_m=\sum_{k=1}^m\frac{\phi^{(k)}(\mu_Y)}{k!}(Y-\mu_Y)^k.
\end{equation}
The constant $\phi(\mu_Y)$ drops out of every Sobol numerator and denominator.
Thus $V_m$ is the reference and $R_m$ is the perturbation.

For $m=1$ with nonzero slope, $V_1$ is affine in $Y$, so the theorem bounds
indices of $\phi(Y)$ relative to indices of $Y$. For $m=2$, the theorem can
handle $Y^2$ and other critical-point transforms whenever the quadratic
reference has positive variance. If the transform is exactly a polynomial and
$m$ includes its degree, then $R_m=0$ and the theorem is exact provided the
polynomial reference is nonconstant. The remaining task is to compute or
estimate the Sobol indices of the polynomial reference.

\subsection*{Why this proof avoids a two-case theorem}

The old split would have required one proof for the nonzero-slope case and a
separate proof for the critical-point case. The projection perturbation theorem
absorbs both. It does not ask whether the first derivative vanishes. It asks only
whether the residual after choosing a reference is small. The two regimes differ
only in the natural choice of reference:
\begin{itemize}
  \item local-affine regime: $V_1=\phi'(\mu_Y)(Y-\mu_Y)$;
  \item critical quadratic regime: $V_2=\frac12\phi''(\mu_Y)(Y-\mu_Y)^2$ when the
    linear term vanishes;
  \item mixed polynomial regime: include all Taylor terms needed to make the
    residual small.
\end{itemize}
The proof is therefore unified, while the interpretation still distinguishes
whether the dominant transformed-output structure is affine, quadratic, or
higher-order.

% ======================================================================
\section{Known limitations and scope for future work}
\label{app:limitations}
% ======================================================================

\paragraph{Reference-relative, not always original-output-relative.}
The main theorem bounds Sobol indices of $\phi(Y)$ relative to the chosen
reference $V_m$. Only when $V_m$ is affine in $Y$ does this become a direct bound
relative to the Sobol indices of $Y$. In critical-point regimes the reference may
be quadratic or higher-order, and its Sobol structure may differ substantially
from that of $Y$. This is the noncommutativity mechanism in the cases that
motivate the theorem.

\paragraph{Choice of Taylor order and computational cost of the reference indices.}
The analyst must choose $m$. A low order gives a simple reference but may leave a
large residual; a higher order may make the residual smaller but produces a more
complicated reference output whose Sobol indices $\mathsf{PE}_{\mathscr C}(V_m)$
appear on the right-hand side of the bound.

This does \emph{not} require additional evaluations of the original simulator.
The reference $V_m$ is a polynomial in the \emph{original} output $Y$
(see \eqref{eq:taylor_reference}), so its Sobol indices can sometimes be
computed analytically from the distributional and ANOVA structure of $Y$, as in
the quadratic example of Section~\ref{sec:numerical}. When an analytic
calculation is not available, $\mathsf{PE}_{\mathscr C}(V_m)$ can be estimated
from the same model evaluations used to estimate the Sobol indices of $Y$, by
post-processing the observed output values into the reference output $V_m$.

If even the reference indices are unavailable, a universal conservative bound
follows by maximizing over $p\in[0,1]$. The right-hand side of
\eqref{eq:abstract_bound} is increasing in $p_{\mathscr C}$ (replacing
$\sqrt{p}$ and $p$ by their maximum value~1 gives an upper bound), yielding
\begin{equation}
  \left|\mathsf{PE}_{\mathscr C}(Z)-\mathsf{PE}_{\mathscr C}(V_m)\right|
  \leq
  \frac{4\eta_m+2\eta_m^2}{(1-\eta_m)^2},
  \label{eq:universal_bound}
\end{equation}
which depends only on $\eta_m$ and requires no knowledge of $\mathsf{PE}_{\mathscr C}(V_m)$.
The theorem is most useful when $\eta_m$ is comfortably below one.

\paragraph{Worst-case conservatism.}
The Cauchy--Schwarz bound is tight only when the relevant ANOVA projections of
the reference and residual are perfectly collinear in $L^2$. In symmetric or
nearly symmetric models the leading residual contribution may be small or zero,
as illustrated by the third-moment contrast \eqref{eq:third_moment_contrast}; in
such cases the observed displacement can be much smaller than the bound.

\paragraph{Bounded support and unbounded extensions.}
The bounded-output assumption in Corollary~\ref{cor:taylor_bound} gives a clean
uniform bound on $\|\phi^{(m+1)}\|_\infty$ and guarantees the required moments.
The abstract Proposition~\ref{prop:abstract_bound} itself does not require
bounded support. The same smooth-transform argument applies in unbounded settings
whenever $V_m\in L^2$, $R_m\in L^2$, $\Var[V_m]>0$, and a finite bound on
$\eta_m$ is available.

\paragraph{Population quantities and estimator uncertainty.}
The inequalities are population statements. Plug-in estimates of
$\eta_m$, $p_{m,\mathscr C}$, $K$, or any Sobol index require separate sampling
uncertainty analysis. Using an observed sample range in place of an almost-sure
support interval gives a diagnostic calculation, not a rigorous population
bound.

\paragraph{Alternative sensitivity targets.}
The theorem concerns standard variance-based Sobol ANOVA component, closed, and
total-effect indices. It does not assert a corresponding perturbation bound for
Borgonovo $\delta$ indices, Cram\'{e}r--von Mises indices, fixed-kernel HSIC
statistics, perturbed-law indices, or other distributional/dependence-based
sensitivity measures. Those methods have different targets and may have
different transformation behavior; the present argument uses the
orthogonal-projection structure of variance-based Sobol decompositions.

\paragraph{Scalar versus spatial or functional aggregates.}
This memo bounds scalar Sobol component indices and the standard scalar closed
and total-effect sums. It does not by itself prove a bound for a
variance-weighted aggregate over spatial, temporal, or functional output
coordinates. Such aggregates use output-coordinate-specific variance weights,
and a nonlinear pointwise transform can change both the local Sobol numerator and
the local variance weight. This is the same distinction emphasized in the
companion memo: pointwise scalar preservation under a common affine transform
immediately implies aggregate preservation, but nonlinear transforms require a
separate aggregate argument. Generalized Sobol indices for Hilbert-space-valued
outputs provide a related but distinct functional-output framework
\citep{Gamboa2014}; for temporal functional outputs, variance-based sensitivity
analysis for time-dependent processes is treated by
\citet{AlexanderianGremaudSmith2020}; spatial and spatiotemporal aggregate
variants require their own weighting and covariance-specific arguments.

\paragraph{Path to a complete article.}
A standalone publication would need: (i) simulation studies comparing the bounds
to empirical shifts across component, closed, and total-effect indices;
(ii) estimator uncertainty quantification for Taylor-reference Sobol indices and
residual ratios; (iii) extensions to unbounded-output settings under moment and
curvature conditions; and (iv) extensions to spatial or functional aggregates.

\bibliographystyle{plainnat}
\bibliography{bounds_memo_v22}

\end{document}
